%% Declaration of generative-AI use. The University of Helsinki requires a student who
%% has used a language model in preparing a thesis to state in writing which model was
%% used and in what way; the model may not be named as an author or cited as a source.
%%
%% Shared verbatim by Thesis_POC.tex and Thesis_v1.tex so the two cannot drift apart.
%% It uses no generated macros and no chapter numbers, because the two documents number
%% their chapters differently.
%%
%% The order of the paragraphs follows the University's basic reporting instruction:
%% what the use covered, what it was for, which tool, how it performed and how its
%% output was checked, and what the author took from it.

\cleardoublepage
\phantomsection
\addcontentsline{toc}{chapter}{Statement on the Use of AI Tools}
\chapter*{Statement on the Use of AI Tools}
\markboth{Statement on the Use of AI Tools}{Statement on the Use of AI Tools}

\noindent\textbf{Tool.} Anthropic's Claude Opus~5, used through the Cursor development
environment, is the only generative AI system used in the preparation of this thesis. It was
used throughout the work, from June to August~2026. Following University of Helsinki policy the
model is neither named as an author nor cited as a source: it cannot take responsibility for
what it produces, and it cannot reliably vouch for the literature it paraphrases.

\noindent\textbf{What the assistance covered.} It was extensive, and is declared as such. The
model drafted and revised the running prose of this thesis and proposed the order in which its
arguments are made. It wrote and refactored the two implementations of the UE Context
Management service, the load generator, the fault-injection catalogue, and the interoperability
prototype. It wrote the analysis code behind the bootstrap confidence intervals, the exact
paired permutation tests, and the Holm correction, together with the figure-generating and
build scripts and the \LaTeX{} sources. It also suggested candidate literature.
% ==================== ACTION REQUIRED BEFORE SUBMISSION ==========================
% The abstract of Thesis_POC.tex was drafted with model assistance. The University's
% thesis guide states that a language model may not be used when preparing the thesis
% summary, and in this programme the summary also serves as the maturity test, where
% AI use is never permitted. Rewrite the abstract unaided before submitting; until
% that is done the closing sentence of this paragraph is not yet true of Thesis_POC.
% =================================================================================
The abstract was written by the author unaided, as the University's thesis guide requires of
the thesis summary.

\noindent\textbf{What the model did not settle.} The research problem, the four research
questions, and the three hypotheses are the author's, as is the control on which the study's
value rests: freezing one language-neutral specification and one concurrency model, so that the
language remains the only systematic difference between the two implementations. So is the
decision that keeps RQ2 falsifiable, namely to carry two defect classes -- a lock-order deadlock
and a lost update -- that lie outside Rust's compile-time guarantee and can therefore embarrass
the premise rather than confirm it. Every measurement reported here was run by the author on the
hosts documented in this thesis; no measured value came from a language model. Each measured
number reaches the text through a macro regenerated from the raw logs, so a figure cannot be
transcribed by hand, by the author or by the model. The readings of those numbers are the
author's as well, including the two restraints the evidence imposes: that a far tail a
virtualised host cannot resolve must not be reported as equivalence, and that a sweep dominated
by cross-thread handoff cannot bound an effect of the language.

\noindent\textbf{How the output was checked.} Model output was treated as a draft to be
falsified rather than as a result. Claims were checked against the sources cited for them: every
bibliography entry was resolved against Crossref, DataCite, or the issuing body, and each was
traced to the page or clause carrying the claim made of it. That audit found four passages where
the prose overstated or misplaced its source, one of them a sentence inside quotation marks that
does not appear in the work it was attributed to; all four were corrected, and the audit is kept
with the artefacts so that the correction remains reviewable. Code was checked by running it:
the two implementations are held to one conformance suite whose transcripts are diffed line by
line, and the fault catalogue fixes the expected outcome of each defect before it is injected,
so that a surprise appears as a discrepancy instead of being explained away afterwards. Nothing
is asserted here on the model's authority.

\noindent\textbf{What the author took from it.} The model is dependable at mechanical work --
transposing an idiom between two languages, restructuring a table, keeping a build honest -- and
undependable in a specific and predictable way: it will supply a plausible citation for a claim,
a plausible cause for a measurement, and a confident sentence where the data supports a hedge.
Its fluency is not evidence. The lasting lesson of the work was accordingly less about either
language than about verification. The effort saved in drafting is spent again, and rightly, in
checking, and the parts of this thesis that are defensible are the parts that were checked
against a source, a specification, or a measurement.

\noindent\textbf{Responsibility.} The author is responsible for the whole of this thesis: its
language, its facts, its code, its measurements, and its conclusions. Where it is wrong, the
error is the author's.
